%% ==========================================================================
%% Conclusion
%% ==========================================================================
\section{Conclusion}\label{sec:conclusion}

Anyone who has used a large language model for research knows the ritual: check
every citation. This paper explains why the ritual exists and why it will not be
engineered away under current architectures.

We have shown that epistemic honesty is unverifiable for systems operating under
text-only observation models with bounded supervision. This is not a capability
limitation to be solved with scale, nor a training failure to be solved with better
RLHF. It is a structural property of the interface: the channel through which
these systems communicate discards the epistemic metadata that would make
verification possible.

The impossibility has a shape. Policies conditioning only on queries cannot
satisfy epistemic honesty across ambiguous world states (\autoref{thm:representational}). Even
grounded architectures cannot learn honest behavior when supervisors cannot
distinguish plausible fabrications from truth (\autoref{thm:learnability}). Stacking judges
does not escape the constraint (\autoref{cor:layered}). Verification cost grows
superlinearly with response complexity (\autoref{lem:superlinear}). Responsibility concentrates
with the system owner (\autoref{cor:responsibility}).

But impossibility results are maps, not walls. The Fischer-Lynch-Paterson
theorem did not end distributed systems; it clarified what assumptions had to
change. Our result plays an analogous role. Systems that export structured
epistemic state—provenance, entropy traces, attention topology—escape the
impossibility class. The tensor interface demonstrates this escape is
achievable with current architectures and minimal overhead.

The tensor interface addresses verification at inference time. A complementary
question is whether training-time interventions can produce systems that are
constitutionally less capable of certain fabrications, such as systems where the
modification class is bounded to exclude dishonest configurations before
deployment. Recent work on pedagogical training suggests bounded modification
during capability development may achieve what personality engineering cannot.
The combination of training-time bounds and inference-time verification would
provide defense in depth: systems less likely to fabricate and more easily
caught when they do.

The work ahead is not exhorting models to try harder at honesty. It is
building systems that make honesty verifiable by design. The tools exist:
vector clocks, cryptographic signatures, content-addressable storage.
The architectural patterns exist. What remains is applying them to the
epistemic dimension with the same rigor we learned to apply to causality
and integrity in distributed systems.

The interface is the problem. The interface is where the solution lives.
